\person{Alda Reed} lifted him then, in the room off \place{the Clay Yard} in \place{Low Kiln}, and the stone did not come with him.

It stayed in the fold of grey wool where his belly had been, a small pale disc with the damp cloth cooling around it, and she did not look at it. She had come to lift a child and clean him, and the lifting was done with the whole weight taken at once, the way a woman lifts a sack of meal, and the cleaning was done with a rag she had warmed by holding it a moment inside her bodice before she touched him. \person{Daniel Marchetti} noted the rag. He noted the temperature of it, which was body-warm and not water-warm, and the order in which she wiped: face, neck, the crease at the groin, the back of the knees, the place where the cord had been cut. It was a good order. In a white room far away he had watched nurses clean a man who could not clean himself and had judged them by just such an order, and he judged her now, and the judgment was useless, because he could not have told her any of it and would not have been believed if he had.

The room went on with its work. \person{Dessa Silt}'s count ran under it, nineteen and eighteen and nineteen, and the fire gave its small irregular tick, and the door that would not shut let the cold in around the edge and went on letting it in. A woman sang somewhere below, and the song had its words worn flat by work, and he had already stopped listening for the tune. \person{Corin} lay where she put him, small and brown and slippery with slip, and did what he was told by the only part of him that had ever learned to obey.

Then the light changed.

Alda turned him to get the back of him, or turned him to wrap him, and the banked fire came up under the ash and put a low red light across the room, and the light went along his ribs and down his left forearm, and there was something on him that had not been there when he was a man.

It was red. Not the red of a bruise coming up under the skin, not the blue-black of a cord or a pressure mark, not the dusted pattern of petechiae he had seen on a dozen faces in the resuscitation bay. The lines ran finer than any of those and they branched. They started somewhere at the third rib on the left side, where the firelight caught them, and went down and outward like a thing growing, and one arm of them crossed the forearm and stopped just short of the wrist, as if it had been told where to stop. The colour was the colour of a coal that has given up its flame and kept its red, a deep ordinary red, not bright. It sat in the skin. There was no swelling at it, no heat he could feel through his own skin, no edge of injury; it was as if somebody had taken a fine brush and a red that came from inside the body and written a small branching sentence across him, and then closed the skin over it.

He made a differential because that was what he did when a thing frightened him. A capillary mark. A bruise older than the bruises of birth. A stain from the cord, or from the hands that had pulled him. A birthmark, which was a word that meant nobody had looked closely enough. He ran through them the way he had run through a hundred rashes and patterns in a white room, and every one of them was wrong in the same way: they were things that happened to the skin, and this was a thing that had been written in it. The lines branched the way a burn branches when it follows warmth along a limb, and there was no burn. He had seen a man who had come out of a fire with just such a map on his arm, the skin whole and the red underneath it, and he had called it a burn map and entered it in a chart. He could enter nothing here. He could not lift his hand to touch the place. He could only lie in the red light and let the mark be the one thing in the room he could not put a name to.

Alda's hands stopped.

It was half a beat, no more --- the pause of a woman who has laid her hands on four hundred children and has never once had reason to take them back --- but Daniel Marchetti had spent his life watching hands stop. He had seen a senior resident's hand stop over a film, and a nurse's hand stop half an inch from a chart, and he had learned to know before anyone spoke that the stop meant the thing on the table had changed and nobody had yet decided what to call it. His own hands, which were soft brown wrists with milky nails, lay where they lay. He could not have pointed at the mark. He could not have said what it was. He knew the shape of the stop, and he filed it, and the filing was the only thing in the room that belonged to him.

Her two fingers left the rag.

They went to the corner of the mark, to the place where the branching ended under his arm, and hovered there over the skin without touching it; then they came back across her own body in a small cutting of the air, forehead and breast and shoulder, finished before the motion had properly begun. It was an old thing done quickly, the way a woman crosses a threshold she has crossed ten thousand times and still touches the wood. It took less than a breath. Then her hands were on the rag again, and the rag was on him, and the moment was gone as completely as if it had never been in the room.

Except that she had seen it. Alda Reed looked up from the child, and her face went through a small change that had nothing in it of surprise and everything of a door being shut quietly. She had made the sign over a child who had come from nowhere anybody could name, and she had made it where the mother might have seen it, and where the aunt standing in the door might have seen it, and there was no account in Grath for a thing like that. She covered it the way a woman covers a spilt cup: with a word.

“\lore[birth-flower]{Birth-flower},” she said.

“What's that?” Dessa said, from the pallet, without opening her eyes.

“That's what they call them.” Alda wiped along his forearm and would not look at the place she had wiped. “They come like that. I saw one in Nine Pots in Y.H. 806, on the fourth day of Stone, on a potter's third child, and there was nothing to it; he's a carter in Stonemouth now.” She blew the dust off the slate without looking at it. “Not a thing to write home about.”

The count stopped. Dessa's eyes opened a little in the tallow light. “Is he hot?” she said.

It was not a question about the mark. It was the question a woman asks about a child who has been given a stone to hold, and Dan knew the shape of it before he knew the words, because he had heard it in a hundred rooms from people standing at the foot of a bed.

“He's warm,” Alda said. “That's the whole of it. Warm and marked and breathing, which is more than some.”

“He's the little hot one,” Dessa said, and the words went flat and finished, and the count went back to nineteen.

Daniel Marchetti heard the date the way he had once heard a man clear his throat before giving a history that was not quite true. She had given a year and a day and a third child and a trade, none of which had been asked for, and all of which had the flat solid sound of a thing placed carefully over a hole. It was the small exactness of a woman choosing what to set down and what to leave out. He had signed a chart once without reading the nurse's note underneath it, and the chart had been exact, and the note had been the thing that mattered, and the institution had called it a systems failure. He recognized the sound of it now in a clay room with a tallow dip, and it made nothing in him louder than curiosity, which was how he knew the infant's body had won again: he wanted the breast, and he wanted to know what the woman had hidden, and the wanting was bigger than the knowing by the whole width of a man's life.

\person{Mira Silt} had not come nearer. She stood just inside the door with the cold still on her sleeves, and her thumb went once into her pocket and came out, and the small hard thing ticked, and she put her hand back on the frame. She had not sat down. She said nothing about the flower. The grey ridge of \lore{the Ash} ran along the skirting board behind her, and nobody in the room would sweep it before the Five.

Then the smell came.

It came sideways, as smells do when a memory has no door to come through: bitter coffee that had sat too long on a hot plate, and under it the clean metal edge of a room scrubbed four times a day with the same mop. The tallow went thin. The clay went thin. The wet wool went away.

The white room is there.

The light in it is overhead and flat, and it comes down on a body on a table and on a nurse holding pressure in a cavity with both hands. Snow comes at the windows in a brown rush. Outside the glass the sodium lamps make orange tunnels in the snow, one tunnel behind another, and the tunnels go up \place{the expressway} toward the split, somewhere in \place{Chicago}, in February, at 2:17 in the morning. A flare stands in the snow and puts a red rim on everything it touches. The flare smokes cheaply; the smoke is acrid in the throat and lies under the sweetness of the fuel.

The wipers have stopped. They have been stopped for some time. Snow builds along the bottom of the windscreen in a fan and the fan grows. The engine is still running. Fuel goes out of the broken van with a hiss that changes over several minutes into something sweeter and worse, and the smell of antifreeze lies under it, and the cold comes in through the door he has left open.

He is working. Airway, bleeding, the end of the problem. There is a woman in the back seat of the minivan who is still moving, and snow has gone up his sleeves, and his fingers have stopped being reliable and he is angry at them. The order of the work does not matter now, because the fuel has caught.

The phone is in his coat pocket against his thigh. It vibrates. The screen is lit where his coat has ridden up, and \person[Nina Marchetti]{a name} is on it, and he does not look at the name. There is a boat on a lake in summer that he has promised for nine years, and a girl who grew up while he worked nights, and a sentence he has composed so often that nothing is left of it but weight. The weight is in the pocket with the phone. It is heavier than the phone. He does not look.

The cold is leaving his hands.

It goes out of them like water out of a wrung cloth, and the heat comes up through his sleeves from every direction at once, a quick enormous heat against a cold that has been there all night, so much of it that it is like being held. There is metal in his mouth. The hiss of the fuel goes on. The orange light goes on. And underneath all of it, finished and complete and with no words in it that anyone could have written down: \emph{I have left the job unfinished.}

At the edge of the flare light \person[Marra]{a woman} stands in a coat the colour of wet ash. She has not come into the light. The flare puts a red rim on the snow at her feet and leaves the rest of her grey, so that she is a shape and nothing more: a coat, a head, a hand. She holds a slate in her hand, and the hand moves. It moves without hurry, short strokes at the top of the slate, then a pause, then more strokes. She has not touched the child. She has not once looked up. The hand moves. It does not stop when the fire takes the van.

Two lights are in the world now. One is small and cold and under his leg, a rectangle of light with a name on it that he will not look at. The other is the flare, red on the snow, and the woman is at the edge of it with the hand moving, and the writing goes on. He holds the two lights together for one breath: the light he refuses, and the light that refuses to stop.

Then the white room was gone.

The smell of coffee stayed in his nose for a breath after the tallow returned, and then that went too. The tallow came back, and the wet wool, and the clay, and the count. Nothing in the room had moved in the time it takes a memory to arrive and go. Dessa was on nineteen. The fire gave its tick. Alda's two fingers came down on the corner of the wool and pressed it flat, the way a woman presses down a cloth that will not stay where it is put, and the word birth-flower settled over the room like a lid.
